% Reviewed translation: PDF pages 315--319 / printed pages 301--305.
\MathBookSourcePage{315}
\subsection{非奇异解分支的逼近}
\label{subsec:navier-nonsingular-branch-approximation}

回到抽象问题 \eqref{eq:navier-branch-abstract-problem} 以引入逼近方法.
对每个趋于零的实参数 $h>0$, 给定从 $\Lambda\times X$ 到 $\mathscr X$ 的
$C^p$ 映射 $F_h$, 它应是 $F$ 的一个逼近. 现在要求解
\begin{equation}\label{eq:navier-branch-discrete-problem}
F_h(\lambda,u_h)=0.
\end{equation}
假设 $\{(\lambda,u(\lambda));\lambda\in\Lambda\}$ 是
\eqref{eq:navier-branch-abstract-problem} 的非奇异解分支. 我们要给出充分条件,
保证 \eqref{eq:navier-branch-discrete-problem} 在该分支的适当邻域中存在唯一解分支.

先固定 $\lambda\in\Lambda$, 并以 \eqref{eq:navier-branch-discrete-problem} 的解
$u_h(\lambda)$ 逼近 $u(\lambda)$. 令 $\widetilde u_h=\widetilde u_h(\lambda)$ 是
$X$ 的任意元素. 引入
\begin{align}
\gamma(\lambda)&=
\|\{D_uF(\lambda,u(\lambda))\}^{-1}\|_{\mathcal L(\mathscr X;X)},
\label{eq:navier-branch-gamma}\\
\mu_h(\lambda)&=
\|D_uF(\lambda,u(\lambda))-D_uF_h(\lambda,\widetilde u_h)\|_{\mathcal L(X;\mathscr X)}.
\label{eq:navier-branch-mu}
\end{align}

\setcounter{statement}{2}
\begin{Lemma}\label{lem:navier-branch-derivative-perturbation}
若
\begin{equation}\label{eq:navier-branch-derivative-smallness}
\gamma(\lambda)\mu_h(\lambda)<1,
\end{equation}
则 $D_uF_h(\lambda,\widetilde u_h)$ 是从 $X$ 到 $\mathscr X$ 的同构.
\end{Lemma}
\begin{Proof}
置
\[
B=\{D_uF(\lambda,u(\lambda))\}^{-1}
\{D_uF(\lambda,u(\lambda))-D_uF_h(\lambda,\widetilde u_h)\}.
\]
则
\[
D_uF_h(\lambda,\widetilde u_h)
=D_uF(\lambda,u(\lambda))(I-B).
\]
由 \eqref{eq:navier-branch-derivative-smallness},
$\|B\|\leq\gamma(\lambda)\mu_h(\lambda)<1$, 因而 $I-B$ 是同构且
\[
\|(I-B)^{-1}\|_{\mathcal L(X;X)}
\leq\frac1{1-\gamma(\lambda)\mu_h(\lambda)}.
\]
所以 $D_uF_h(\lambda,\widetilde u_h)$ 也是同构, 并且
\begin{equation}\label{eq:navier-branch-inverse-perturbation-bound}
\|\{D_uF_h(\lambda,\widetilde u_h)\}^{-1}\|_{\mathcal L(\mathscr X;X)}
\leq\frac{\gamma(\lambda)}{1-\gamma(\lambda)\mu_h(\lambda)}.
\end{equation}
\end{Proof}

\setcounter{statement}{0}
\begin{Remark}\label{rem:navier-branch-common-range-space}
上述 Lemma 利用了 $F$ 和 $F_h$ 都把 $\Lambda\times X$ 映入 $\mathscr X$ 这一事实.
这一假设使后面的理论比 $F_h$ 与 $F$ 不定义在同一空间时更简单, 但它并非下面不动点论证的基本工具.
Subsection~\ref{subsec:navier-nondifferentiable-branches} 将给出二者不定义在同一空间的例子.
\end{Remark}

\MathBookSourcePage{316}
现在假设 $D_uF_h(\lambda,\widetilde u_h)$ 是从 $X$ 到 $\mathscr X$ 的同构, 并记
\begin{equation}\label{eq:navier-branch-local-quantities}
\left\{
\begin{aligned}
\varepsilon_h(\lambda)&=\|F_h(\lambda,\widetilde u_h)\|_{\mathscr X},\\
\gamma_h(\lambda)&=\|\{D_uF_h(\lambda,\widetilde u_h)\}^{-1}\|_{\mathcal L(\mathscr X;X)},\\
S(u;\alpha)&=\{v\in X;\ \|v-u\|_X\leq\alpha\},\\
L_h(\lambda;\alpha)&=
\sup_{v\in S(\widetilde u_h;\alpha)}
\|D_uF_h(\lambda,\widetilde u_h)-D_uF_h(\lambda,v)\|_{\mathcal L(X;\mathscr X)}.
\end{aligned}
\right.
\end{equation}
显然 $\alpha\mapsto L_h(\lambda;\alpha)$ 从 $\mathbb R_+$ 到 $\mathbb R_+$ 单调递增.

\setcounter{statement}{0}
\begin{Theorem}\label{thm:navier-branch-local-existence-error}
假设
\begin{align}
D_uF_h(\lambda,\widetilde u_h)&
\text{ is an isomorphism of }X\text{ onto }\mathscr X,
\label{eq:navier-branch-local-isomorphism}\\
2\gamma_h(\lambda)L_h(\lambda;2\gamma_h(\lambda)\varepsilon_h(\lambda))&<1.
\label{eq:navier-branch-local-contraction}
\end{align}
则 \eqref{eq:navier-branch-discrete-problem} 有一个解 $(\lambda,u_h)$, 且
\begin{equation}\label{eq:navier-branch-local-ball}
u_h\in S(\widetilde u_h;2\gamma_h(\lambda)\varepsilon_h(\lambda)).
\end{equation}
此外 $D_uF_h(\lambda,u_h)$ 是从 $X$ 到 $\mathscr X$ 的同构, 且
\begin{equation}\label{eq:navier-branch-solution-inverse-bound}
\|\{D_uF_h(\lambda,u_h)\}^{-1}\|_{\mathcal L(\mathscr X;X)}
\leq2\gamma_h(\lambda).
\end{equation}
再者, 对每个满足
\begin{equation}\label{eq:navier-branch-uniqueness-radius}
\gamma_h(\lambda)L_h(\lambda;\alpha)<1
\end{equation}
的半径 $\alpha$, $u_h$ 是球 $S(\widetilde u_h;\alpha)$ 中的唯一解, 并有估计
\begin{equation}\label{eq:navier-branch-fundamental-error}
\|u_h-v\|_X
\leq\frac{\gamma_h(\lambda)}
{1-\gamma_h(\lambda)L_h(\lambda;\alpha)}
\|F_h(\lambda,v)\|_{\mathscr X}
\quad\forall v\in S_h(u_h;\alpha).
\end{equation}
\end{Theorem}
\begin{Proof}
定义
\[
\Phi_h(v)=v-[D_uF_h(\lambda,\widetilde u_h)]^{-1}F_h(\lambda,v).
\]
$(\lambda,u_h)$ 是 \eqref{eq:navier-branch-discrete-problem} 的解当且仅当 $u_h$ 是
$\Phi_h$ 的不动点. 置
$S=S(\widetilde u_h;2\gamma_h(\lambda)\varepsilon_h(\lambda))$.

对 $v\in S$, Taylor 公式给出
\[
F_h(\lambda,v)-F_h(\lambda,\widetilde u_h)
=\int_0^1D_uF_h(\lambda,\widetilde u_h+\theta(v-\widetilde u_h))
\cdot(v-\widetilde u_h)\,d\theta.
\]
因而由 \eqref{eq:navier-branch-local-quantities} 和
\eqref{eq:navier-branch-local-contraction},
\[
\|\Phi_h(v)-\widetilde u_h\|_X
\leq\gamma_h\varepsilon_h
\{2\gamma_hL_h(\lambda;2\gamma_h\varepsilon_h)+1\}
<2\gamma_h\varepsilon_h.
\]
所以 $\Phi_h(S)\subset S$.

\MathBookSourcePage{317}
同理, 对 $v,w\in S$,
\[
\|\Phi_h(v)-\Phi_h(w)\|_X
\leq\gamma_h(\lambda)L_h(\lambda;2\gamma_h\varepsilon_h)
\|v-w\|_X<\frac12\|v-w\|_X.
\]
故 $\Phi_h$ 是 $S$ 到自身的严格压缩映射, 并在 $S$ 中有唯一不动点 $u_h$.
又有
\[
\|D_uF_h(\lambda,u_h)-D_uF_h(\lambda,\widetilde u_h)\|
\leq L_h(\lambda;2\gamma_h\varepsilon_h)<\frac1{2\gamma_h}.
\]
由 Lemma~\ref{lem:navier-branch-derivative-perturbation} 和
\eqref{eq:navier-branch-inverse-perturbation-bound} 得
\eqref{eq:navier-branch-solution-inverse-bound}.

若 $v$ 是 $\Phi_h$ 在 $S(\widetilde u_h;\alpha)$ 中的另一不动点, 则
\[
\|v-u_h\|_X\leq\gamma_h(\lambda)L_h(\lambda;\alpha)\|v-u_h\|_X,
\]
由 \eqref{eq:navier-branch-uniqueness-radius} 得 $v=u_h$.
最后, 对任意 $v\in S(\widetilde u_h;\alpha)$,
\begin{equation}\label{eq:navier-branch-error-identity}
\begin{aligned}
v-u_h=[D_uF_h(\lambda,\widetilde u_h)]^{-1}
\bigg[&\int_0^1\{D_uF_h(\lambda,\widetilde u_h)
-D_uF_h(\lambda,u_h+\theta(v-u_h))\}\cdot(v-u_h)\,d\theta\\
&+F_h(\lambda,v)\bigg].
\end{aligned}
\end{equation}
因此
\[
\|v-u_h\|_X
\leq\gamma_h(\lambda)
\{L_h(\lambda;\alpha)\|v-u_h\|_X+\|F_h(\lambda,v)\|_{\mathscr X}\},
\]
这就证明了 \eqref{eq:navier-branch-fundamental-error}.
\end{Proof}

\MathBookSourcePage{318}
现在允许 $\lambda$ 在紧区间 $\Lambda$ 中变化, 并用 $C^0$ 映射
$\lambda\mapsto\widetilde u_h(\lambda):\Lambda\to X$ 代替固定元素.

\setcounter{statement}{3}
\begin{Lemma}\label{lem:navier-branch-uniform-isomorphism}
若
\begin{equation}\label{eq:navier-branch-uniform-mu}
\lim_{h\to0}\sup_{\lambda\in\Lambda}\mu_h(\lambda)=0,
\end{equation}
则存在 $h_0>0$, 使对所有 $\lambda\in\Lambda$ 和 $h\leq h_0$,
$D_uF_h(\lambda,\widetilde u_h(\lambda))$ 是从 $X$ 到 $\mathscr X$ 的同构, 且
\begin{equation}\label{eq:navier-branch-uniform-gamma}
\gamma_h(\lambda)=
\|[D_uF_h(\lambda,\widetilde u_h(\lambda))]^{-1}\|_{\mathcal L(\mathscr X;X)}
\leq2\gamma(\lambda).
\end{equation}
\end{Lemma}
\begin{Proof}
因 $\Lambda$ 紧, $\gamma(\lambda)$ 在 $\Lambda$ 上有界.
由 \eqref{eq:navier-branch-uniform-mu}, 对某个 $h_0>0$,
\[
\sup_{\lambda\in\Lambda}\gamma(\lambda)\mu_h(\lambda)\leq\frac12
\qquad\forall h\leq h_0.
\]
应用 Lemma~\ref{lem:navier-branch-derivative-perturbation} 和
\eqref{eq:navier-branch-inverse-perturbation-bound} 即得结论.
\end{Proof}

\setcounter{statement}{1}
\begin{Theorem}\label{thm:navier-branch-uniform-branch}
设 $\{(\lambda,u(\lambda));\lambda\in\Lambda\}$ 是
\eqref{eq:navier-branch-abstract-problem} 的非奇异解分支, 且给定
$\lambda\mapsto\widetilde u_h(\lambda)\in C^0(\Lambda;X)$ 满足
\eqref{eq:navier-branch-uniform-mu}. 若还满足
\begin{align}
\lim_{h\to0}\sup_{\lambda\in\Lambda}\varepsilon_h(\lambda)&=0,
\label{eq:navier-branch-uniform-residual}\\
\lim_{\alpha\to0}\left\{\sup_{\lambda\in\Lambda}L_h(\lambda;\alpha)\right\}&=0
\quad\text{uniformly for all }h\leq h_0,
\label{eq:navier-branch-uniform-lipschitz}
\end{align}
则存在 $\alpha>0$, $h_1>0$ 以及 $u_h\in C^0(\Lambda;X)$, 使对所有
$h\leq h_1$:
\begin{equation}\label{eq:navier-branch-uniform-nonsingular-branch}
\{(\lambda,u_h(\lambda));\lambda\in\Lambda\}
\text{ is a branch of nonsingular solutions of }
\eqref{eq:navier-branch-discrete-problem},
\end{equation}
并且对每个 $\lambda\in\Lambda$, $u_h(\lambda)$ 是球
$S(\widetilde u_h(\lambda);\alpha)$ 中的唯一解, 且
\begin{equation}\label{eq:navier-branch-uniform-error}
\|u_h(\lambda)-v\|_X
\leq4\gamma(\lambda)\|F_h(\lambda,v)\|_{\mathscr X}
\quad\forall v\in S(\widetilde u_h(\lambda);\alpha).
\end{equation}
\end{Theorem}
\begin{Proof}
由 Lemma~\ref{lem:navier-branch-uniform-isomorphism}, $\gamma_h$ 对
$\lambda\in\Lambda$, $h\leq h_0$ 一致有界. 由
\eqref{eq:navier-branch-uniform-lipschitz}, 可取 $\alpha>0$ 使
\begin{equation}\label{eq:navier-branch-uniform-half}
\gamma_h(\lambda)L_h(\lambda;\alpha)<\frac12
\quad\forall\lambda\in\Lambda,\quad\forall h\leq h_0.
\end{equation}
又由 \eqref{eq:navier-branch-uniform-residual}, 存在 $h_0'>0$ 使
\[
2\gamma_h(\lambda)L_h(\lambda;2\gamma_h(\lambda)\varepsilon_h(\lambda))<1
\quad\forall\lambda\in\Lambda,\quad\forall h\leq h_0'.
\]
取 $h_1=\min(h_0,h_0')$. Theorem~\ref{thm:navier-branch-local-existence-error}
给出每个 $\lambda$ 的唯一非奇异解 $u_h(\lambda)$.
再由 Lemma~\ref{lem:navier-branch-uniform-isomorphism} 和
\eqref{eq:navier-branch-solution-inverse-bound},
\begin{equation}\label{eq:navier-branch-uniform-solution-inverse}
\|[D_uF_h(\lambda,u_h(\lambda))]^{-1}\|_{\mathcal L(\mathscr X;X)}
\leq4\gamma(\lambda).
\end{equation}
公式 \eqref{eq:navier-branch-fundamental-error},
\eqref{eq:navier-branch-uniform-gamma} 和
\eqref{eq:navier-branch-uniform-half} 给出
\eqref{eq:navier-branch-uniform-error}. 最后, $u_h$ 关于 $\lambda$ 的连续性由
\eqref{eq:navier-branch-uniform-error}、$\widetilde u_h$ 的连续性以及 $F_h$ 的连续性直接得到.
\end{Proof}

\MathBookSourcePage{319}
\setcounter{statement}{1}
\begin{Remark}\label{rem:navier-branch-arbitrary-center}
上述证明的一个重要特点是 $\widetilde u_h$ 可以任意选择. 当 $F$ 和 $F_h$
定义在同一空间 $X$ 上时, 最直接的选择是
$\widetilde u_h(\lambda)=u(\lambda)$. 当 $F_h$ 定义在不同于 $X$ 的空间
$X_h$ 上时, 如 Subsection~\ref{subsec:navier-nondifferentiable-branches},
则取 $u$ 的适当逼近.
\end{Remark}

\setcounter{statement}{2}
\begin{Remark}\label{rem:navier-branch-regularity-in-parameter}
Theorem~\ref{thm:navier-branch-uniform-branch} 没有充分使用 $F$ 和 $F_h$
关于 $\lambda$ 的正则性. 若 $\lambda\mapsto\widetilde u_h(\lambda)$ 是从
$\Lambda$ 到 $X$ 的 $C^p$ 映射, 则 Crouzeix~\cite{CrouzeixForthcoming}
还证明了 $u_h\in C^p(\Lambda;X)$.
\end{Remark}
